\IfFileExists{stacks-project.cls}{%
\documentclass{stacks-project}
}{%
\documentclass{amsart}
}

% For dealing with references we use the comment environment
\usepackage{verbatim}
\newenvironment{reference}{\comment}{\endcomment}
%\newenvironment{reference}{}{}
\newenvironment{slogan}{\comment}{\endcomment}
\newenvironment{history}{\comment}{\endcomment}

% For commutative diagrams we use Xy-pic
\usepackage[all]{xy}

% We use 2cell for 2-commutative diagrams.
\xyoption{2cell}
\UseAllTwocells

% We use multicol for the list of chapters between chapters
\usepackage{multicol}

% This is generally recommended for better output
\usepackage{lmodern}
\usepackage[T1]{fontenc}

% For cross-file-references
\usepackage{xr-hyper}

% Package for hypertext links:
\usepackage{hyperref}

% For any local file, say "hello.tex" you want to link to please
% use \externaldocument[hello-]{hello}
\externaldocument[introduction-]{introduction}
\externaldocument[conventions-]{conventions}
\externaldocument[sets-]{sets}
\externaldocument[categories-]{categories}
\externaldocument[topology-]{topology}
\externaldocument[sheaves-]{sheaves}
\externaldocument[sites-]{sites}
\externaldocument[stacks-]{stacks}
\externaldocument[fields-]{fields}
\externaldocument[algebra-]{algebra}
\externaldocument[brauer-]{brauer}
\externaldocument[homology-]{homology}
\externaldocument[derived-]{derived}
\externaldocument[simplicial-]{simplicial}
\externaldocument[more-algebra-]{more-algebra}
\externaldocument[smoothing-]{smoothing}
\externaldocument[modules-]{modules}
\externaldocument[sites-modules-]{sites-modules}
\externaldocument[injectives-]{injectives}
\externaldocument[cohomology-]{cohomology}
\externaldocument[sites-cohomology-]{sites-cohomology}
\externaldocument[dga-]{dga}
\externaldocument[dpa-]{dpa}
\externaldocument[sdga-]{sdga}
\externaldocument[hypercovering-]{hypercovering}
\externaldocument[schemes-]{schemes}
\externaldocument[constructions-]{constructions}
\externaldocument[properties-]{properties}
\externaldocument[morphisms-]{morphisms}
\externaldocument[coherent-]{coherent}
\externaldocument[divisors-]{divisors}
\externaldocument[limits-]{limits}
\externaldocument[varieties-]{varieties}
\externaldocument[topologies-]{topologies}
\externaldocument[descent-]{descent}
\externaldocument[perfect-]{perfect}
\externaldocument[more-morphisms-]{more-morphisms}
\externaldocument[flat-]{flat}
\externaldocument[groupoids-]{groupoids}
\externaldocument[more-groupoids-]{more-groupoids}
\externaldocument[etale-]{etale}
\externaldocument[chow-]{chow}
\externaldocument[intersection-]{intersection}
\externaldocument[pic-]{pic}
\externaldocument[weil-]{weil}
\externaldocument[adequate-]{adequate}
\externaldocument[dualizing-]{dualizing}
\externaldocument[duality-]{duality}
\externaldocument[discriminant-]{discriminant}
\externaldocument[derham-]{derham}
\externaldocument[local-cohomology-]{local-cohomology}
\externaldocument[algebraization-]{algebraization}
\externaldocument[curves-]{curves}
\externaldocument[resolve-]{resolve}
\externaldocument[models-]{models}
\externaldocument[functors-]{functors}
\externaldocument[equiv-]{equiv}
\externaldocument[pione-]{pione}
\externaldocument[etale-cohomology-]{etale-cohomology}
\externaldocument[proetale-]{proetale}
\externaldocument[relative-cycles-]{relative-cycles}
\externaldocument[more-etale-]{more-etale}
\externaldocument[trace-]{trace}
\externaldocument[crystalline-]{crystalline}
\externaldocument[spaces-]{spaces}
\externaldocument[spaces-properties-]{spaces-properties}
\externaldocument[spaces-morphisms-]{spaces-morphisms}
\externaldocument[decent-spaces-]{decent-spaces}
\externaldocument[spaces-cohomology-]{spaces-cohomology}
\externaldocument[spaces-limits-]{spaces-limits}
\externaldocument[spaces-divisors-]{spaces-divisors}
\externaldocument[spaces-over-fields-]{spaces-over-fields}
\externaldocument[spaces-topologies-]{spaces-topologies}
\externaldocument[spaces-descent-]{spaces-descent}
\externaldocument[spaces-perfect-]{spaces-perfect}
\externaldocument[spaces-more-morphisms-]{spaces-more-morphisms}
\externaldocument[spaces-flat-]{spaces-flat}
\externaldocument[spaces-groupoids-]{spaces-groupoids}
\externaldocument[spaces-more-groupoids-]{spaces-more-groupoids}
\externaldocument[bootstrap-]{bootstrap}
\externaldocument[spaces-pushouts-]{spaces-pushouts}
\externaldocument[spaces-chow-]{spaces-chow}
\externaldocument[groupoids-quotients-]{groupoids-quotients}
\externaldocument[spaces-more-cohomology-]{spaces-more-cohomology}
\externaldocument[spaces-simplicial-]{spaces-simplicial}
\externaldocument[spaces-duality-]{spaces-duality}
\externaldocument[formal-spaces-]{formal-spaces}
\externaldocument[restricted-]{restricted}
\externaldocument[spaces-resolve-]{spaces-resolve}
\externaldocument[formal-defos-]{formal-defos}
\externaldocument[defos-]{defos}
\externaldocument[cotangent-]{cotangent}
\externaldocument[examples-defos-]{examples-defos}
\externaldocument[algebraic-]{algebraic}
\externaldocument[examples-stacks-]{examples-stacks}
\externaldocument[stacks-sheaves-]{stacks-sheaves}
\externaldocument[criteria-]{criteria}
\externaldocument[artin-]{artin}
\externaldocument[quot-]{quot}
\externaldocument[stacks-properties-]{stacks-properties}
\externaldocument[stacks-morphisms-]{stacks-morphisms}
\externaldocument[stacks-limits-]{stacks-limits}
\externaldocument[stacks-cohomology-]{stacks-cohomology}
\externaldocument[stacks-perfect-]{stacks-perfect}
\externaldocument[stacks-introduction-]{stacks-introduction}
\externaldocument[stacks-more-morphisms-]{stacks-more-morphisms}
\externaldocument[stacks-geometry-]{stacks-geometry}
\externaldocument[moduli-]{moduli}
\externaldocument[moduli-curves-]{moduli-curves}
\externaldocument[examples-]{examples}
\externaldocument[exercises-]{exercises}
\externaldocument[guide-]{guide}
\externaldocument[desirables-]{desirables}
\externaldocument[coding-]{coding}
\externaldocument[obsolete-]{obsolete}
\externaldocument[fdl-]{fdl}
\externaldocument[index-]{index}

% Theorem environments.
%
\theoremstyle{plain}
\newtheorem{theorem}[subsection]{Theorem}
\newtheorem{proposition}[subsection]{Proposition}
\newtheorem{lemma}[subsection]{Lemma}

\theoremstyle{definition}
\newtheorem{definition}[subsection]{Definition}
\newtheorem{example}[subsection]{Example}
\newtheorem{exercise}[subsection]{Exercise}
\newtheorem{situation}[subsection]{Situation}

\theoremstyle{remark}
\newtheorem{remark}[subsection]{Remark}
\newtheorem{remarks}[subsection]{Remarks}

\numberwithin{equation}{subsection}

% Macros
%
\def\lim{\mathop{\mathrm{lim}}\nolimits}
\def\colim{\mathop{\mathrm{colim}}\nolimits}
\def\Spec{\mathop{\mathrm{Spec}}}
\def\Hom{\mathop{\mathrm{Hom}}\nolimits}
\def\Ext{\mathop{\mathrm{Ext}}\nolimits}
\def\SheafHom{\mathop{\mathcal{H}\!\mathit{om}}\nolimits}
\def\SheafExt{\mathop{\mathcal{E}\!\mathit{xt}}\nolimits}
\def\Sch{\mathit{Sch}}
\def\Mor{\mathop{\mathrm{Mor}}\nolimits}
\def\Ob{\mathop{\mathrm{Ob}}\nolimits}
\def\Sh{\mathop{\mathit{Sh}}\nolimits}
\def\NL{\mathop{N\!L}\nolimits}
\def\CH{\mathop{\mathrm{CH}}\nolimits}
\def\proetale{{pro\text{-}\acute{e}tale}}
\def\etale{{\acute{e}tale}}
\def\QCoh{\mathit{QCoh}}
\def\Ker{\mathop{\mathrm{Ker}}}
\def\Im{\mathop{\mathrm{Im}}}
\def\Coker{\mathop{\mathrm{Coker}}}
\def\Coim{\mathop{\mathrm{Coim}}}

% Boxtimes
%
\DeclareMathSymbol{\boxtimes}{\mathbin}{AMSa}{"02}

%
% Macros for moduli stacks/spaces
%
\def\QCohstack{\mathcal{QC}\!\mathit{oh}}
\def\Cohstack{\mathcal{C}\!\mathit{oh}}
\def\Spacesstack{\mathcal{S}\!\mathit{paces}}
\def\Quotfunctor{\mathrm{Quot}}
\def\Hilbfunctor{\mathrm{Hilb}}
\def\Curvesstack{\mathcal{C}\!\mathit{urves}}
\def\Polarizedstack{\mathcal{P}\!\mathit{olarized}}
\def\Complexesstack{\mathcal{C}\!\mathit{omplexes}}
% \Pic is the operator that assigns to X its picard group, usage \Pic(X)
% \Picardstack_{X/B} denotes the Picard stack of X over B
% \Picardfunctor_{X/B} denotes the Picard functor of X over B
\def\Pic{\mathop{\mathrm{Pic}}\nolimits}
\def\Picardstack{\mathcal{P}\!\mathit{ic}}
\def\Picardfunctor{\mathrm{Pic}}
\def\Deformationcategory{\mathcal{D}\!\mathit{ef}}
\usepackage{cleveref}
\crefname{originalequation}{original equation}{original equations}
\Crefname{originalequation}{Original equation}{Original equations}
\makeatletter
\newlabel{equation-valuative}{{0H1V}{source}{}{sourcepointer-1}{}}
\newlabel{equation-valuative@cref}{{[originalequation][0][]0H1V}{[1][0][]source}{}{}{}}
\makeatother
\crefname{originallemma}{original lemma}{original lemmas}
\Crefname{originallemma}{Original lemma}{Original lemmas}
\makeatletter
\newlabel{lemma-Noetherian-dvr-valuative-separation}{{0H1W}{source}{}{sourcepointer-2}{}}
\newlabel{lemma-Noetherian-dvr-valuative-separation@cref}{{[originallemma][0][]0H1W}{[1][0][]source}{}{}{}}
\makeatother
\crefname{originalreference}{original reference}{original references}
\Crefname{originalreference}{Original reference}{Original references}
\makeatletter
\newlabel{limits-lemma-Noetherian-dvr-valuative-proper}{{0H1X}{source}{}{sourcepointer-3}{}}
\newlabel{limits-lemma-Noetherian-dvr-valuative-proper@cref}{{[originalreference][0][]0H1X}{[1][0][]source}{}{}{}}
\makeatother
\crefname{originalreference}{original reference}{original references}
\Crefname{originalreference}{Original reference}{Original references}
\makeatletter
\newlabel{spaces-cohomology-lemma-weak-chow}{{089J}{source}{}{sourcepointer-4}{}}
\newlabel{spaces-cohomology-lemma-weak-chow@cref}{{[originalreference][0][]089J}{[1][0][]source}{}{}{}}
\makeatother
\crefname{originalreference}{original reference}{original references}
\Crefname{originalreference}{Original reference}{Original references}
\makeatletter
\newlabel{spaces-morphisms-definition-valuative-criterion}{{03IX}{source}{}{sourcepointer-5}{}}
\newlabel{spaces-morphisms-definition-valuative-criterion@cref}{{[originalreference][0][]03IX}{[1][0][]source}{}{}{}}
\makeatother
\crefname{originalreference}{original reference}{original references}
\Crefname{originalreference}{Original reference}{Original references}
\makeatletter
\newlabel{spaces-morphisms-lemma-characterize-proper}{{0A40}{source}{}{sourcepointer-6}{}}
\newlabel{spaces-morphisms-lemma-characterize-proper@cref}{{[originalreference][0][]0A40}{[1][0][]source}{}{}{}}
\makeatother
\crefname{originalreference}{original reference}{original references}
\Crefname{originalreference}{Original reference}{Original references}
\makeatletter
\newlabel{spaces-morphisms-lemma-image-proper-is-proper}{{08AJ}{source}{}{sourcepointer-7}{}}
\newlabel{spaces-morphisms-lemma-image-proper-is-proper@cref}{{[originalreference][0][]08AJ}{[1][0][]source}{}{}{}}
\makeatother
\crefname{originalreference}{original reference}{original references}
\Crefname{originalreference}{Original reference}{Original references}
\makeatletter
\newlabel{spaces-morphisms-lemma-universally-closed-local}{{03IT}{source}{}{sourcepointer-8}{}}
\newlabel{spaces-morphisms-lemma-universally-closed-local@cref}{{[originalreference][0][]03IT}{[1][0][]source}{}{}{}}
\makeatother
\crefname{originalreference}{original reference}{original references}
\Crefname{originalreference}{Original reference}{Original references}
\makeatletter
\newlabel{spaces-properties-proposition-locally-quasi-separated-open-dense-scheme}{{06NH}{source}{}{sourcepointer-9}{}}
\newlabel{spaces-properties-proposition-locally-quasi-separated-open-dense-scheme@cref}{{[originalreference][0][]06NH}{[1][0][]source}{}{}{}}
\makeatother
\title{Properness from the dvr valuative criterion over a Noetherian base}
\author{The Stacks Project Authors}
\date{stacks-project:spaces-limits.tex:local-snapshot-20261012; GNU FDL 1.2 or later}
\begin{document}
\maketitle
\noindent\textit{Editorial note.} Counted span is the author's own proof environment, source lines 4918-4984, following the statement at 4896-4916.\par\medskip
\noindent\textit{Editorial note.} The single '(details omitted)' covers replacing X' by the disjoint union of the irreducible components dominating those of X, flagged in the notes as standard bookkeeping.\par\medskip
\noindent\textit{Editorial note.} The footnote on the sharper formulation (existence after an extension of discrete valuation rings) is transcribed; it does not change the equivalence.\par\medskip
\noindent\textit{Editorial note.} Steps the author leaves to the reader are flagged here for the reviewer: source line(s) 4962 carry an "omitted"/"follows easily"/"immediate from the definitions" phrase. Those sentences are retained verbatim; the displayed argument around them stays inside the counted span.\par\medskip
\noindent\textit{Allowed standard prerequisites (editorial).}\begin{itemize}
\item Characterization of proper morphisms of spaces. Morphisms of Spaces, Lemma spaces-morphisms-lemma-characterize-proper (cross-chapter, cited at \textasciitilde{}4919-4920). Equivalences (1) iff (2) iff (3).
\item Separation from the dvr valuative criterion. lemma-Noetherian-dvr-valuative-separation (same chapter, cited at \textasciitilde{}4925). f is separated; reduces to universally closed.
\item Universally closed is etale-local on the base. Morphisms of Spaces, Lemma spaces-morphisms-lemma-universally-closed-local (cross-chapter, cited at \textasciitilde{}4928-4931). Reduces to the affine etale base change V x\_Y X -> V.
\item Weak Chow's lemma for algebraic spaces. Cohomology of Spaces, Lemma spaces-cohomology-lemma-weak-chow (cross-chapter, cited at \textasciitilde{}4958-4960). Quasi-projective X' over Y with a proper surjective X' -> X.
\item Image of a proper morphism is proper. Morphisms of Spaces, Lemma spaces-morphisms-lemma-image-proper-is-proper (cross-chapter, cited at \textasciitilde{}4966-4967). X is proper over Y once X' -> Y is proper.
\item Scheme-level dvr valuative criterion. Limits, Lemma limits-lemma-Noetherian-dvr-valuative-proper (cross-chapter, cited at \textasciitilde{}4968-4969). Criterion for X' -> Y proper in terms of generic residue fields.
\end{itemize}\medskip
\ifdefined\hypertarget\hypertarget{sourcepointer-1}{}\fi
\noindent\textit{Source reference.} Original equation 0H1V: sources/stacks-spaces-limits/pinned-original-source.tex, source lines 4818--4818. Navigation pointer only: the counted proof cites this statement by \textbackslash{}ref as a standing input; the deduction the counted proof has to make is carried out inside the counted span, and no part of that deduction is delegated to this pointer.\par\medskip
\ifdefined\hypertarget\hypertarget{sourcepointer-2}{}\fi
\noindent\textit{Source reference.} Original lemma 0H1W: sources/stacks-spaces-limits/pinned-original-source.tex, source lines 4829--4829. Navigation pointer only: the counted proof cites this statement by \textbackslash{}ref as a standing input; the deduction the counted proof has to make is carried out inside the counted span, and no part of that deduction is delegated to this pointer.\par\medskip
\ifdefined\hypertarget\hypertarget{sourcepointer-3}{}\fi
\noindent\textit{Source reference.} Original reference 0H1X: focused-seeds/geometry/S9-spaces-limits-Noetherian-dvr-valuative-proper/ref-tags.txt, source lines 5--5. Navigation pointer only: the counted proof cites this statement by \textbackslash{}ref as a standing input; the deduction the counted proof has to make is carried out inside the counted span, and no part of that deduction is delegated to this pointer.\par\medskip
\ifdefined\hypertarget\hypertarget{sourcepointer-4}{}\fi
\noindent\textit{Source reference.} Original reference 089J: focused-seeds/geometry/S9-spaces-limits-Noetherian-dvr-valuative-proper/ref-tags.txt, source lines 6--6. Navigation pointer only: the counted proof cites this statement by \textbackslash{}ref as a standing input; the deduction the counted proof has to make is carried out inside the counted span, and no part of that deduction is delegated to this pointer.\par\medskip
\ifdefined\hypertarget\hypertarget{sourcepointer-5}{}\fi
\noindent\textit{Source reference.} Original reference 03IX: focused-seeds/geometry/S9-spaces-limits-Noetherian-dvr-valuative-proper/ref-tags.txt, source lines 7--7. Navigation pointer only: the counted proof cites this statement by \textbackslash{}ref as a standing input; the deduction the counted proof has to make is carried out inside the counted span, and no part of that deduction is delegated to this pointer.\par\medskip
\ifdefined\hypertarget\hypertarget{sourcepointer-6}{}\fi
\noindent\textit{Source reference.} Original reference 0A40: focused-seeds/geometry/S9-spaces-limits-Noetherian-dvr-valuative-proper/ref-tags.txt, source lines 8--8. Navigation pointer only: the counted proof cites this statement by \textbackslash{}ref as a standing input; the deduction the counted proof has to make is carried out inside the counted span, and no part of that deduction is delegated to this pointer.\par\medskip
\ifdefined\hypertarget\hypertarget{sourcepointer-7}{}\fi
\noindent\textit{Source reference.} Original reference 08AJ: focused-seeds/geometry/S9-spaces-limits-Noetherian-dvr-valuative-proper/ref-tags.txt, source lines 9--9. Navigation pointer only: the counted proof cites this statement by \textbackslash{}ref as a standing input; the deduction the counted proof has to make is carried out inside the counted span, and no part of that deduction is delegated to this pointer.\par\medskip
\ifdefined\hypertarget\hypertarget{sourcepointer-8}{}\fi
\noindent\textit{Source reference.} Original reference 03IT: focused-seeds/geometry/S9-spaces-limits-Noetherian-dvr-valuative-proper/ref-tags.txt, source lines 10--10. Navigation pointer only: the counted proof cites this statement by \textbackslash{}ref as a standing input; the deduction the counted proof has to make is carried out inside the counted span, and no part of that deduction is delegated to this pointer.\par\medskip
\ifdefined\hypertarget\hypertarget{sourcepointer-9}{}\fi
\noindent\textit{Source reference.} Original reference 06NH: focused-seeds/geometry/S9-spaces-limits-Noetherian-dvr-valuative-proper/ref-tags.txt, source lines 11--11. Navigation pointer only: the counted proof cites this statement by \textbackslash{}ref as a standing input; the deduction the counted proof has to make is carried out inside the counted span, and no part of that deduction is delegated to this pointer.\par\medskip
\begin{lemma}
\label{lemma-Noetherian-dvr-valuative-proper}
Let $S$ be a scheme. Let $f : X \to Y$ be a morphism of algebraic spaces
over $S$. Assume $f$ is quasi-separated and of finite type and $Y$
is locally Noetherian. The following are equivalent:
\begin{enumerate}
\item $f$ is proper,
\item $f$ satisfies the valuative criterion, see
Morphisms of Spaces, Definition
\ref{spaces-morphisms-definition-valuative-criterion},
\item for any diagram (\ref{equation-valuative}) there exists exactly
one dotted arrow,
\item for all diagrams (\ref{equation-valuative}) with $A$ a discrete
valuation ring there exists exactly one dotted arrow, and
\item for all diagrams (\ref{equation-valuative}) where $A$ is a discrete
valuation ring and where the image of $\Spec(K) \to X$ is a point of
codimension $0$ on $X$ there exists exactly one dotted arrow\footnote{There is
a sharper formulation where in the existence part one only requires
the dotted arrow exists after an extension of discrete valuation rings.}.
\end{enumerate}
\end{lemma}

\begin{proof}
We have (1) $\Leftrightarrow$ (2) $\Leftrightarrow$ (3) by 
Morphisms of Spaces, Lemma \ref{spaces-morphisms-lemma-characterize-proper}.
It is clear that (3) $\Rightarrow$ (4) $\Rightarrow$ (5).
To finish the proof we will now show (5) implies (1).

\medskip\noindent
Assume (5). By Lemma \ref{lemma-Noetherian-dvr-valuative-separation}
we see that $f$ is separated. To finish the proof it suffices to show
that $f$ is universally closed. Let $V \to Y$ be an \'etale morphism
where $V$ is an affine scheme. It suffices to show that the base
change $V \times_Y X \to V$ is universally closed, see
Morphisms of Spaces, Lemma
\ref{spaces-morphisms-lemma-universally-closed-local}.
Let
$$
\xymatrix{
\Spec(K) \ar[r] \ar[d] & V \times_Y X \ar[d] \ar[r] & X \ar[d] \\
\Spec(A) \ar[r] \ar@{-->}[ru] \ar@{..>}[rru] & V \ar[r] & Y
}
$$
of algebraic spaces over $S$
be a commutative diagram where $A$ is a discrete valuation ring
with fraction field $K$ and where $\Spec(K) \to V \times_Y X$
maps to a point of codimension $0$ of the algebraic space
$V \times_Y X$. Since $V \times_Y X \to X$ is \'etale it follows
that the image of $\Spec(K) \to X$ is a point of codimension $0$
of $X$. Thus by (5) we obtain the longer of the two dotted
arrows fitting into the diagram. Then of course we obtain the
shorter one as well. It follows that our assumptions hold for
the morphism $V \times_Y X \to V$ and we reduce to the case discussed
in the next paragraph.

\medskip\noindent
Aassume $Y$ is a Noetherian affine scheme. In this case $X$ is a
separated Noetherian algebraic space (we already know $f$ is separated)
of finite type over $Y$. (In particular, the algebraic space $X$
has a dense open subspace which is a scheme by Properties of Spaces, Proposition
\ref{spaces-properties-proposition-locally-quasi-separated-open-dense-scheme}
although strictly speaking we will not need this.)
Choose a quasi-projective scheme $X'$ over $Y$ and a proper surjective
morphism $X' \to X$ as in the weak form of Chow's lemma
(Cohomology of Spaces, Lemma \ref{spaces-cohomology-lemma-weak-chow}).
We may replace $X'$ by the disjoint union of the irreducible components
which dominate an irreducible component of $X$; details omitted.
In particular, we may assume that generic points of the scheme $X'$ map to
points of codimension $0$ of $X$ (in this case these are exactly
the generic points of $X$). We claim that $X' \to Y$ is proper.
The claim implies $X$ is proper over $Y$ by
Morphisms of Spaces, Lemma \ref{spaces-morphisms-lemma-image-proper-is-proper}.
To prove this, according to
Limits, Lemma \ref{limits-lemma-Noetherian-dvr-valuative-proper}
it suffices to prove that in every solid commutative diagram
$$
\xymatrix{
\Spec(K) \ar[r] \ar[d] & X' \ar[r] & X \ar[d] \\
\Spec(A) \ar[rr] \ar@{-->}[ru]^a \ar@{-->}[rru]_b & & Y
}
$$
where $A$ is a dvr with fraction field $K$ and where $K$ is the residue
field of a generic point of $X'$ we can find the
dotted arrow $a$ (we already know uniqueness as $X'$ is separated).
By assumption (5) we can find the dotted arrow $b$.
Then the morphism $X' \times_{X, b} \Spec(A) \to \Spec(A)$
is a proper morphism of schemes and by the valuative criterion
for morphisms of schemes we can lift $b$ to the desired morphism $a$.
\end{proof}

\end{document}
